\subsection{Datasets}

\paragraph{MSLesSeg.}
The primary experiments use the MSLesSeg
dataset~\cite{guarnera_mslesseg_2025} of longitudinal brain MRI scans of
patients with multiple sclerosis. Of its 75 subjects, the 53 with manual
lesion annotations are used here, as the official test set is unlabeled.
Each subject has one to four timepoints, each with T1-weighted (T1-w),
T2-weighted (T2-w), and Fluid-Attenuated Inversion Recovery (FLAIR)
images. The dataset organizers provide the images skull-stripped,
co-registered to the MNI152 template, and resampled to 1~mm isotropic
spacing.

\paragraph{3D-MR-MS.}
To assess whether the behavior of CATMIL carries over to other data, the
four losses are also trained and evaluated on a second multiple sclerosis
dataset, 3D-MR-MS~\cite{lesjak_novel_2018}, which was acquired at a
different center with different imaging conditions. It contains 30
patients with consensus-based white matter lesion annotations, provided
co-registered and bias-field corrected. Of the available modalities, only
T1-w, T2-w, and FLAIR are used, to keep the input configuration identical
to MSLesSeg. The voxels of this dataset are anisotropic
($0.47 \times 0.47 \times 0.8$~mm), so lesion sizes on 3D-MR-MS are given
in mm$^3$ rather than in voxels.

\subsection{Data Split}

Data are partitioned at the patient level, so that scans of one patient
never appear in both training and evaluation sets. For MSLesSeg, six
patients, with all their timepoints (12 test cases), are held out as an
independent test set; for 3D-MR-MS, six of the 30 patients are held out.
The remaining patients (47 and 24, respectively) are divided into five
folds following the nnU-Net cross-validation strategy. Each of the five
fold models is evaluated on the test set, and results are reported as the
mean and sample standard deviation across the five models; statistical
comparisons are described in \cref{sec:stats}. All training
settings, including the CATMIL loss weights, are identical for the two
datasets, so 3D-MR-MS was not used to tune the objective.

\subsection{Compared Losses}

Four configurations of nnU-Net are compared, differing only in the loss
function:
\begin{itemize}
    \item \textbf{DiceCE}: the default Dice and cross-entropy loss;
    \item \textbf{Tversky}~\cite{salehi_tversky_2017}: Tversky loss, with
    $\alpha=0.3$ and $\beta=0.7$ weighting false positives and false
    negatives;
    \item \textbf{FocalTversky}~\cite{abraham_novel_2019}: Focal Tversky
    loss, $\mathcal{L}_{\mathrm{Tversky}}^{1/\gamma_{\mathrm{F}}}$, with the
    same $\alpha$ and $\beta$ and focal exponent $\gamma_{\mathrm{F}}=4/3$;
    \item \textbf{CATMIL}: the proposed loss.
\end{itemize}
Tversky and FocalTversky replace the DiceCE loss, whereas CATMIL adds its
two terms to it. Ablations of the two CATMIL terms and of their weights are
reported in \cref{appendix:catmil-components,appendix:catmil-loss-weights}.

\subsection{Implementation}

All models use nnU-Net (version 2.1.1)~\cite{isensee_nnu-net_2021} with
its default preprocessing, patch-based training, and optimizer: stochastic
gradient descent with Nesterov momentum of 0.99, weight decay of
$3\times10^{-5}$, and an initial learning rate of 0.01 with polynomial
decay. Training runs for 150 epochs with a batch size of 2 on an NVIDIA
Quadro RTX 8000 GPU. The loss implementation and evaluation code are
available at \url{https://github.com/luumsk/SmallLesionMRI}.

\subsection{Evaluation Metrics}

Three groups of metrics are used: voxel-level metrics for overall
segmentation quality, lesion-level metrics for detection, in particular of
small lesions, and error-volume metrics for the extent of false-positive
and false-negative errors. Lesion-level metrics are included because the
DSC is dominated by large lesions (\cref{eq:sl_missed_lesion}) and does not
reflect whether individual, clinically relevant lesions are
found~\cite{basit_rethinking_2026,jaus_every_2025}.

\paragraph{Voxel-level metrics.}
Overall segmentation quality is measured with the DSC~\cite{dice_measures_1945}
and the 95th-percentile Hausdorff distance (HD95, in
mm)~\cite{taha_metrics_2015}. The name Dice is kept only for the Dice loss
and the DiceCE baseline, which use the differentiable form of the
coefficient.

\paragraph{Lesion definition and detection.}
A lesion is a 3D connected component of the binary mask under
6-connectivity, the same definition as in the proposed loss. A reference
lesion counts as detected if at least one predicted voxel overlaps it, and
a predicted component counts as correct if it overlaps at least one
reference lesion. This any-overlap criterion is lenient: a prediction that
covers only part of a lesion still counts as a detection. To test whether
the results depend on it, lesion recall was also computed with stricter
criteria that require at least 10\%, 25\%, or 50\% of the lesion voxels to
be covered (\cref{appendix:strict-detection}).

\paragraph{Lesion-level metrics.}
Lesion-wise recall is the fraction of reference lesions that are detected,
lesion-wise precision is the fraction of predicted components that are
correct, and the lesion-wise F1 score combines the two. Because precision
counts components regardless of their size, many small false-positive
components lower it substantially even when their total volume is small.
Small Lesion Recall is the lesion-wise recall restricted to reference
lesions of at most 150 voxels, which corresponds to 150~mm$^3$ at the 1~mm
isotropic spacing of MSLesSeg and to about 26~mm$^3$ on 3D-MR-MS. The FN
Count is the number of reference lesions that are completely missed. The
recall of lesions of at least 3~mm is the lesion-wise recall restricted to
lesions above the size used in clinical MRI reading: current consensus
recommendations state that multiple sclerosis lesions should generally be
at least 3~mm in longest diameter, with exceptions in specific locations
such as the corpus callosum and the floor of the fourth
ventricle~\cite{barkhof_2024_2025,filippi_mri_2016}. The longest diameter
of a reference lesion is measured as its largest in-plane extent over all
axial, coronal, and sagittal slices.

\paragraph{Error-volume metrics.}
The FP Volume is the physical volume of the false-positive voxels. The FP
Blob Fraction is the share of this volume that lies in false-positive
components more than 2~mm from any reference lesion, separating detached
false alarms from extensions of true lesions. The FN Volume Fraction is
the share of the total reference lesion volume that belongs to missed
lesions.

All metrics are computed per test case and averaged over cases; their
exact definitions are given in \cref{appendix:metrics}.

\subsection{Statistical Analysis}
\label{sec:stats}

The standard deviations in the result tables describe the variation
between the five fold models. They do not describe the uncertainty that
comes from the small test sets, so the differences between CATMIL and each
baseline are also tested with the patient as the independent unit. This
matters for MSLesSeg, whose 12 test cases are longitudinal scans of six
patients.

For each metric, the difference between CATMIL and a baseline is computed
in the same way as the table values: the mean over test cases of the
five-model mean. Its 95\% confidence interval (CI) is obtained with a
two-level bootstrap with 10{,}000 resamples. Each resample draws patients
with replacement, together with all their scans, and draws fold models with
replacement. The fold models are drawn jointly for both losses, because
fold $k$ of every loss is trained on the same data split. For lesion-level
recall, the same bootstrap is applied to the pooled reference lesions of
the resampled patients. A paired exact Wilcoxon signed-rank test is applied
to the per-patient differences, with the Holm correction over the three
baselines.

With six patients, the smallest possible two-sided $p$-value of the exact
test is $1/32 \approx 0.031$ before correction and $0.094$ after it, so no
comparison can reach $p < 0.05$ after correction. A difference is
therefore regarded as supported when its 95\% CI excludes zero, and the
number of patients in which CATMIL is better is reported with it. These
analyses describe the consistency of an effect across the available
patients; they are not population-level estimates.
